\documentclass[11pt,a4paper]{article}
\usepackage[T1]{fontenc}
\usepackage[utf8]{inputenc}
\usepackage{lmodern}
\usepackage[margin=27mm]{geometry}
\usepackage{amsmath,amssymb,amsthm,microtype,enumitem,xcolor,fancyhdr}
\usepackage[hidelinks]{hyperref}
\definecolor{ink}{HTML}{183449}
\definecolor{muted}{HTML}{536477}
\hypersetup{pdftitle={Efficient learning from disjunction agreement pairs},
 pdfauthor={Grisha Pochuev}}
\setlength{\parskip}{5pt}
\setlength{\parindent}{0pt}
\setlist{itemsep=2pt,topsep=4pt}
\pagestyle{fancy}
\fancyhf{}
\fancyhead[L]{\small\color{muted}Disjunction agreement pairs}
\fancyhead[R]{\small\color{muted}20 October 2026}
\fancyfoot[C]{\small\thepage}
\setlength{\headheight}{14pt}
\newtheorem{theorem}{Theorem}
\newtheorem{lemma}[theorem]{Lemma}
\newcommand{\C}{\mathcal C}
\newcommand{\HU}{\mathcal H_U}
\newcommand{\E}{\mathbb E}
\newcommand{\Prb}{\mathbb P}
\newcommand{\CONS}{\operatorname{CONS}}

\begin{document}
\begin{center}
{\LARGE\bfseries\color{ink}Efficient learning from\\[3pt]
disjunction agreement pairs}\\[9pt]
{\large An algorithm and proof for the formulation in Blum--Balcan (2007)}\\[10pt]
Grisha Pochuev\\
\href{mailto:n_854@mail.ru}{\texttt{n\_854@mail.ru}}\\
20 October 2026
\end{center}

\begin{abstract}
An elementary exclusion-closure algorithm computes the greatest disjunction
consistent with labelled examples and exact agreement on unlabelled pairs.
Combining this solver with two finite-class bounds gives a polynomial-time
learner using
\(O((\ln|\C_{D,\chi}(\tau)|+\ln(1/\delta))/\varepsilon)\)
labelled pairs and
\(O((n+\ln(1/\delta))/\tau)\) unlabelled pairs, for arbitrary joint
distributions satisfying the realizable agreement promise. Taking
\(\tau=\varepsilon\) gives the labelled-sample benchmark in the stated
2007 question. A fully specified stopping rule removes the need to know
or count the compatible class, with the same order of labels in expectation
and with high probability.
\end{abstract}

\begin{quote}
\small\textbf{Status.} This note supplies a self-contained mathematical argument
for the stated learning guarantee. It does not establish that the argument is
new, that the problem remained open until this date, or that the historical
reward is currently payable. Those are separate bibliographic and author-review
questions. No external acceptance is claimed.
\end{quote}

\section{Setting and precise guarantee}

Blum and Balcan's 2007 note \cite{bb07}, Section 3, asks for efficient
learning with a number of labelled examples controlled by the logarithm
of the nearly compatible class. The relevant item in Section 3.2 is
\emph{Co-training with disjunctions}. The reward statement in that note
offers \$300 for a positive solution to an item in the list.

We state the result using \(d\) fixed Boolean features
\(\phi_1,\ldots,\phi_d\) on an object space \(X\). The hypothesis class consists
of the distinct functions
\[
 \C=\left\{h_S:S\subseteq[d]\right\},\qquad
 h_S(x)=\bigvee_{i\in S}\phi_i(x).
\]
The empty disjunction is \(0\). Repeated representations of the same Boolean
function are counted only once; thus \(|\C|\le 2^d\). Ordinary monotone
disjunctions have \(d=n\), with \(\phi_i(x)=x_i\).
Allowing negated literals is handled in Section~\ref{sec:literals}.

Let \(D\) be any distribution on \(X\times X\). There is a fixed
\(h^*=h_{S^*}\in\C\) such that
\[
 \Prb_{(x_1,x_2)\sim D}\bigl[h^*(x_1)=h^*(x_2)\bigr]=1.
\]
A labelled pair has the common label \(y=h^*(x_1)=h^*(x_2)\).
All sampled pairs are independent draws from \(D\); the coordinates
inside a pair can be arbitrarily dependent and can have different marginals.

Define the incompatibility probability and a prediction error by
\begin{align}
 q_D(h)&=\Prb_D[h(x_1)\ne h(x_2)],\\
 e_D(h)&=\Prb_D[h(x_1)\ne h^*(x_1)
                 \ \text{or}\ h(x_2)\ne h^*(x_2)].
\end{align}
The second quantity is stronger than error on the first view, which is the
prediction convention in \cite{bb07}. Either first-view or second-view error
is at most \(e_D(h)\). For \(0<\tau<1\), set
\[
 \C_D(\tau)=\{h\in\C:q_D(h)\le\tau\},\qquad K=|\C_D(\tau)|.
\]
This is \(\C_{D,\chi}(\tau)\) for the compatibility in the question.
Since \(h^*\) is perfectly compatible, \(1\le K\le 2^d\).

\begin{theorem}[Fixed sample sizes]\label{thm:main}
Let \(0<\varepsilon,\tau<1\) and \(0<\delta\le1/2\).
Draw \(u\) unlabelled pairs and an independent set of \(m\) labelled pairs, where
\begin{equation}\label{eq:sizes}
 u\ \ge\ \frac{d\ln2+\ln(2/\delta)}{\tau},
 \qquad
 m\ \ge\ \frac{\ln K+\ln(2/\delta)}{\varepsilon}.
\end{equation}
The exclusion-closure algorithm below returns a disjunction \(\widehat h\)
satisfying, with probability at least \(1-\delta\),
\[
 e_D(\widehat h)\le\varepsilon,\qquad q_D(\widehat h)\le\tau.
\]
It runs in time polynomial in \(d,u,m\).
For \(d\ge1\), with dense Boolean input, a queue implementation takes \(O(d(u+m)+d)\)
elementary operations after the features have been evaluated.

The algorithm itself takes the samples as input and does not use \(K\).
Section~\ref{sec:adaptive} gives a complete rule for acquiring and stopping
the labelled sample without knowing \(K\).
\end{theorem}

\section{An exact consistency solver}

Write \(A(x)=\{i:\phi_i(x)=1\}\). For an unlabelled pair, write
\(A=A(x_1)\), \(B=A(x_2)\). If \(E=[d]\setminus S\) denotes excluded
features, then
\begin{equation}\label{eq:closed}
 h_S(x_1)=h_S(x_2)
 \quad\Longleftrightarrow\quad
 (A\subseteq E\ \Longleftrightarrow\ B\subseteq E).
\end{equation}
Thus each pair gives two implications between sets:
\begin{equation}\label{eq:rules}
 A\subseteq E\ \Longrightarrow\ B\subseteq E,
 \qquad
 B\subseteq E\ \Longrightarrow\ A\subseteq E.
\end{equation}
An implication fires only when \emph{every} feature in its premise has
already been excluded. An overlap with the excluded set does not suffice.

\medskip
\textbf{Algorithm \(\CONS(U,L)\).}
\begin{enumerate}
\item Let \(E\) initially contain every feature active in either coordinate
of any negatively labelled pair in \(L\).
\item Apply all rules \eqref{eq:rules} from \(U\), adding their conclusions
to \(E\), until no further feature can be added.
\item Set \(\widehat S=[d]\setminus E\) and return \(h_{\widehat S}\).
For a general consistency-decision interface, check all positive labelled
pairs and report inconsistency if one is not satisfied.
Under the realizability promise this check always succeeds.
\end{enumerate}

\begin{lemma}[Safety of exclusion]\label{lem:safety}
Throughout the algorithm, \(E\cap S^*=\varnothing\).
Consequently \(S^*\subseteq\widehat S\), and the returned disjunction has
no false negatives anywhere on \(X\).
\end{lemma}

\begin{proof}
If a pair has label \(0\), neither of its coordinates can contain an active
feature from \(S^*\). Initialization therefore excludes no target feature.

Suppose the invariant holds before a rule fires, and suppose \(A\subseteq E\).
Then \(A\cap S^*=\varnothing\), so \(h^*(x_1)=0\). Compatibility of the
target on this sampled pair implies \(h^*(x_2)=0\), hence
\(B\cap S^*=\varnothing\). Adding \(B\) to \(E\) preserves the invariant.
The argument for the other direction is identical.

All sampled pairs satisfy the target promise almost surely, including in the
finite capped procedure used later. Finally \(S^*\subseteq\widehat S\)
implies \(h^*(x)\le h_{\widehat S}(x)\) for every \(x\in X\).
\end{proof}

\begin{lemma}[Consistency and maximality]\label{lem:cons}
The algorithm terminates, agrees on every pair in \(U\), and satisfies every
labelled pair in \(L\). Its support is the greatest feasible support, not just
a maximal one. On arbitrary labelled input, failure of its positive-label
check proves that no feasible disjunction exists.
\end{lemma}

\begin{proof}
Each nontrivial update adds at least one feature to \(E\), so at most \(d\)
such updates occur. At termination both implications in \eqref{eq:rules}
hold for every pair, proving empirical agreement.

Every feature in a negatively labelled pair was excluded initially, so both
coordinates receive label \(0\). If a pair has label \(1\), each coordinate
has an active target feature. By Lemma~\ref{lem:safety} these features are
retained, so both coordinates receive label \(1\).

For maximality, replace \(S^*\) in the safety proof by any support \(T\)
that agrees on \(U\) and avoids every negative labelled example.
The same induction shows \(T\subseteq\widehat S\).
The returned support itself has those properties. It is therefore their
greatest feasible support. Positive constraints are monotone in the support:
if any feasible support satisfies a positive example, so does
\(\widehat S\). If \(\widehat S\) fails a positive constraint, every smaller
feasible support fails it as well.
\end{proof}

\subsection{Why the computation is efficient}

A simple implementation scans all pairs after each change and already
takes polynomial time, at most \(O(d^2u+dm+d)\) on dense input.
A queue implementation is linear in the total input incidence size:
\begin{enumerate}
\item For each of the \(2u\) pair sides, store the number of its features not
yet excluded.
\item For each feature, store the pair sides in which it occurs.
\item When a feature is excluded for the first time, decrement all its
incidence counters. Enqueue a side when its counter reaches zero.
\item On removing a zero-count side from the queue, exclude all features
on the opposite side of that pair.
\end{enumerate}
Each feature is excluded at most once; each side reaches zero at most once;
and each occurrence is used in at most one decrement and one opposite-side
scan. All sides whose counters are initially zero, including empty sides,
are enqueued. Preprocessing and checking
labelled input are also linear in its incidence size.

Equivalently, putting \(z_i=1\) when feature \(i\) is excluded, a pair
contributes Horn implications
\[
 \bigwedge_{i\in A}z_i\ \Longrightarrow\ z_j \quad(j\in B),
 \qquad
 \bigwedge_{j\in B}z_j\ \Longrightarrow\ z_i \quad(i\in A).
\]
The queue groups the common premises, avoiding an unnecessary expansion
into \(|A||B|\) repeated literal occurrences. No satisfiability or counting
oracle is required.

\section{Proof of the labelled-sample guarantee}

For the unlabelled sample \(U\), define the empirical compatible class
\[
 \HU=\{h\in\C:h(x_1)=h(x_2)\text{ for every }(x_1,x_2)\in U\}.
\]
This class is only used in the proof; the learner never enumerates it.

\begin{lemma}[Unlabelled reduction]\label{lem:reduction}
For every \(u\ge0\),
\[
 \Prb[\HU\not\subseteq\C_D(\tau)]\le|\C|e^{-\tau u}.
\]
\end{lemma}
\begin{proof}
For a fixed \(h\) with \(q_D(h)>\tau\), the probability of agreeing on all
\(u\) independent pairs is
\[
 (1-q_D(h))^u\le e^{-q_D(h)u}\le e^{-\tau u}.
\]
Sum this bound over the at most \(|\C|\) hypotheses outside \(\C_D(\tau)\).
\end{proof}

\begin{lemma}[Conditional labelled bound]\label{lem:label}
For a fixed \(U\), if \(L\) consists of \(m\) independent labelled pairs
drawn independently of \(U\), then
\[
 \Prb[\exists h\in\HU:e_D(h)>\varepsilon
            \text{ and }h\text{ satisfies }L\mid U]
 \le|\HU|e^{-\varepsilon m}.
\]
\end{lemma}
\begin{proof}
The class \(\HU\) is fixed after conditioning on \(U\).
A fixed hypothesis of pair prediction error \(p\) satisfies an independent
labelled pair on both coordinates with probability \(1-p\).
Its probability of satisfying all \(m\) pairs is \((1-p)^m\le e^{-pm}\).
Sum over \(\HU\). No independence of the two coordinates is used.
\end{proof}

\begin{proof}[Proof of Theorem~\ref{thm:main}]
By the first inequality in \eqref{eq:sizes} and Lemma~\ref{lem:reduction},
with probability at least \(1-\delta/2\) the event
\[
 \mathcal E_U=\{\HU\subseteq\C_D(\tau)\}
\]
holds. On that event \(|\HU|\le K\), and every possible algorithmic output
has incompatibility at most \(\tau\).

The solver returns a member of \(\HU\) satisfying all labels, by
Lemma~\ref{lem:cons}. Conditional on any such \(U\),
Lemma~\ref{lem:label} and the second inequality in \eqref{eq:sizes} give
\[
 \Prb[e_D(\widehat h)>\varepsilon\mid U]
 \le K e^{-\varepsilon m}\le\delta/2.
\]
Combining the two failure probabilities proves both asserted error bounds.
The algorithmic analysis proves the running-time claim.
\end{proof}

In particular, for \(\tau=\varepsilon\) and \(d=n\) or \(2n\),
\[
 \boxed{\quad
 u=O\!\left(\frac{n+\ln(1/\delta)}{\varepsilon}\right),\qquad
 m=O\!\left(\frac{\ln|\C_{D,\chi}(\varepsilon)|
                         +\ln(1/\delta)}{\varepsilon}\right).
 \quad}
\]
The use of unlabelled data is substantive: the logarithm in the labelled
bound is that of the nearly compatible class, not the full \(2^d\) class.

\section{Stopping without knowing the compatible class}\label{sec:adaptive}

The fixed-sample theorem does not require counting \(K\) to run the solver.
To remove even the need to choose a label budget using an unknown \(K\),
the following fully specified wrapper uses fresh random labels and
independent validation. It does not make active label queries or select
examples for labelling.

Let
\begin{equation}\label{eq:adaptive}
 N=\left\lceil\frac{d\ln2+\ln(4/\delta)}{\tau}\right\rceil,\quad
 b=\ln(16/\delta),\quad
 m_j=\left\lceil\frac{64\,2^j b}{\varepsilon}\right\rceil\quad(j\ge0).
\end{equation}
Draw \(N\) unlabelled pairs once. The computable final phase is
\[
 J=\min\left\{j\ge0:
             m_j\ge\frac4\varepsilon(d\ln2+b)\right\}.
\]
The constant \(64\) is deliberately conservative; \(J\) is allowed to be \(0\).

For \(j=0,\ldots,J\):
\begin{enumerate}
\item Draw \(m_j\) fresh labelled training pairs and compute \(h_j=\CONS(U,L_j)\).
\item If \(j=J\), return \(h_j\).
\item Otherwise draw \(m_j\) fresh labelled validation pairs.
Return \(h_j\) if the fraction on which at least one view is misclassified
is at most \(\varepsilon/2\). If not, continue.
\end{enumerate}

\begin{theorem}[Automatic label budget]\label{thm:adaptive}
The wrapper always terminates in polynomial time. With probability at
least \(1-\delta\), it simultaneously returns a hypothesis with
\[
 e_D(\widehat h)\le\varepsilon,\qquad q_D(\widehat h)\le\tau
\]
and uses at most
\begin{equation}\label{eq:400}
 Q\le400\,\frac{\ln K+\ln(16/\delta)}{\varepsilon}
\end{equation}
labelled pairs, counting both training and validation.
Moreover,
\[
 \E Q=O\!\left(\frac{\ln K+\ln(1/\delta)}{\varepsilon}\right).
\]
A deterministic worst-case bound uses \(d\ln2\) instead of \(\ln K\).
The bound through \(K\) is a high-probability and an expectation guarantee.
\end{theorem}

\subsection{Accuracy and high-probability label count}

The event \(\mathcal E_U\) now fails with probability at most \(\delta/4\).
Conditional on a current hypothesis and all prior information, its
validation block remains independent. Standard binomial bounds give:
\[
 \begin{array}{ll}
 e_D(h_j)>\varepsilon
 &\Longrightarrow\
 \Prb[\text{validation accepts}]\le e^{-\varepsilon m_j/8},\\[3pt]
 e_D(h_j)\le\varepsilon/4
 &\Longrightarrow\
 \Prb[\text{validation rejects}]\le e^{-\varepsilon m_j/12}.
 \end{array}
\]
The first is the lower-tail bound at half the mean. For the second,
binomial stochastic domination reduces to mean \(\varepsilon/4\), and
the upper-tail bound is applied at twice that mean.

The probability of ever accepting an error greater than \(\varepsilon\)
before the final phase is at most
\[
 \sum_{j\ge0}e^{-\varepsilon m_j/8}
 \le\sum_{j\ge0}e^{-8\,2^j b}
 \le2e^{-8b}\le\delta/16.
\]
At phase \(J\), Lemma~\ref{lem:label} over the entire class gives
\[
 \Prb[e_D(h_J)>\varepsilon/4\mid U,\text{past}]
 \le 2^d e^{-\varepsilon m_J/4}
 \le e^{-b}=\delta/16.
\]
These bounds do not require \(\mathcal E_U\).

For analysis only, define the unknown phase
\[
 j^*=\min\left\{j\ge0:
       m_j\ge\frac4\varepsilon(\ln K+b)\right\}.
\]
Since \(K\le2^d\), \(j^*\le J\).
Conditional on \(\mathcal E_U\), if phase \(j^*\) is reached, its fresh
training sample yields \(e_D(h_{j^*})\le\varepsilon/4\) except with
probability \(e^{-b}=\delta/16\). If \(j^*<J\), rejection of such a
hypothesis has probability at most \(e^{-\varepsilon m_{j^*}/12}\le
e^{-b}=\delta/16\). If \(j^*=J\), it stops regardless.

Therefore the probability, on \(\mathcal E_U\), of continuing beyond
\(j^*\) is at most \(\delta/8\). Conditioning on histories is legitimate
because each training and validation block is fresh; alternatively,
all independent blocks can be generated in advance.

On stopping by \(j^*\), the number of labels obeys
\[
 Q\le2\sum_{j=0}^{j^*}m_j\le6m_{j^*}
 \le384\,\frac{\ln K+b}{\varepsilon}+6
 \le400\,\frac{\ln K+b}{\varepsilon}.
\]
For the middle inequalities, use \(\lceil c2^j\rceil\le c2^j+1\),
\(j^*+1\le2^{j^*}\), and minimality of \(j^*\).
If \(j^*=0\), \(m_0\le64b/\varepsilon+1\);
otherwise \(m_{j^*}<8(\ln K+b)/\varepsilon+1\).

The total failure probability for the two errors and the label bound is
at most
\[
 \delta/4+\delta/16+\delta/16+\delta/8=\delta/2\le\delta.
\]
Replacing \(K\) by \(2^d\) in the same size calculation bounds the
deterministic cap. In particular,
\(J=O(\log(1+d/\ln(1/\delta)))\), and all time and sample bounds are
polynomial in the ordinary learning parameters.

\subsection{Expected label count}

For a fixed \(U\), put \(k_U=|\HU|\ge1\) and define \(j_U\) by replacing
\(K\) with \(k_U\) in the definition of \(j^*\). At every reached phase
\(j\ge j_U\), the conditional probability of continuing is at most
\[
 r=2e^{-b}=\delta/8\le1/8.
\]
This follows from the same training and validation bounds; at the final
phase continuation probability is zero. Also \(m_{j+1}\le2m_j\).
Consequently the expected cost after phase \(j_U\) is at most
\[
 \sum_{s\ge1}2^{s+1}m_{j_U}r^s
 =2m_{j_U}\frac{2r}{1-2r}\le\frac23m_{j_U}.
\]
Together with the cost of earlier phases,
\[
 \E[Q\mid U]\le7m_{j_U}
 \le448\,\frac{\ln k_U+b}{\varepsilon}+7.
\]

It remains to control rare, unhelpful unlabelled samples without paying
an extra term proportional to \(\delta d\).
Every \(h\notin\C_D(\tau)\) belongs to \(\HU\) with probability at most
\(e^{-\tau N}\). Hence
\[
 \E k_U\le K+|\C|e^{-\tau N}\le K+\delta/4.
\]
Jensen's inequality gives
\[
 \E\ln k_U\le\ln(K+\delta/4)\le\ln K+\delta/4.
\]
Substitution proves the stated expectation bound. This also shows why no
class-counting procedure is hidden inside the stopping rule.

\section{Literal conventions and extensions}\label{sec:literals}

\subsection{Negated literals}

For ordinary Boolean disjunctions with negations, use the feature map
\[
 \phi(x)=(x_1,\ldots,x_n,1-x_1,\ldots,1-x_n),
 \qquad d=2n.
\]
The consistency proof applies to this restricted set of feature vectors
without change. It never assumes that all \(2^{2n}\) feature vectors can
occur.

The counting in \(K\) is over distinct Boolean functions. In particular,
every literal subset containing both \(x_i\) and \(\neg x_i\) represents
the same constant-one function. One can canonicalize all such outputs to
\(1\). The full signed disjunction class, including the empty disjunction,
has at most \(3^n+1\) distinct functions; the looser bound \(2^{2n}\) is
sufficient for the unlabelled sample.

If a convention excludes constant functions, adjoining them changes
\(K\) by at most two, so the order of the labelled bound is unchanged.
In that optional convention the output may be outside the restricted
class. The principal theorem uses the class of all disjunctions generated
by the stated literals, including the empty disjunction and tautologies
when representable.

\subsection{Different disjunctions in two views}

The same argument also covers targets \(h_1^*,h_2^*\) on two different
feature sets, promised to agree on the pair distribution.
Use disjoint feature indices for the two views, take the two active sets
in that common index universe, and apply the same two implications.
The compatible-class count then refers to distinct hypothesis pairs.
No conditional-independence assumption is needed for this extension.

\section{Small example and verification}

Consider five features and the unlabelled pairs
\[
 (\{2\},\{3,4\}),\qquad
 (\{4\},\{5\}),\qquad
 (\{1,3\},\{1,5\}).
\]
Suppose a labelled negative pair is \((\{2\},\{2\})\).
Initially \(E=\{2\}\).
The first agreement pair forces \(3,4\) into \(E\);
the second forces \(5\) into \(E\).
The third pair imposes no further exclusion. The output is
\(\widehat h(x)=x_1\). A positive labelled pair
\((\{1,3\},\{1,5\})\) is satisfied automatically.

The accompanying code implements the queue solver, signed feature lifting,
and the fully specified wrapper. An independently written finite enumerator
checked the underlying closure property on 19,560 systems:
5,560 exhaustive two-pair systems with all exclusion masks for \(d\le3\);
10,000 random systems for \(4\le d\le8\); and 4,000 target-generated
labelled systems. It checked 48,000 labelled pairs in the last group.
No counterexample was found.

A separate verifier compared the actual queue implementation against
complete enumeration on 22,560 systems, including 2,000 signed-literal
systems and 1,000 arbitrary, potentially inconsistent systems.
All 160,627 feasible-support witnesses and all 71,335 deletion-certificate
steps passed. It correctly rejected 746 inconsistent systems.
Exact results are in \texttt{solver\_verification\_results.json}.

These are finite consistency checks, not a machine-checked formalization
of the probabilistic theorems. The mathematical proofs above establish
the all-dimensions guarantee.

\section{Bibliographic and reward status}

The 2013 paper \cite{bbel13} studies two-sided disjunctions, where two
OR functions give complementary labels to a \emph{single} object, with
additional distributional conditions in its passive results.
The 2024 paper \cite{dmrt24} studies co-training under a weak-dependence
condition using active learning. Neither set of hypotheses and resource
bounds directly supplies the arbitrary-distribution, passive guarantee
proved here.

The later author version of the framework paper \cite{bb10}, Section 6.2,
does not repeat the particular co-training-disjunction item, though it
does list other related questions. This omission does not establish
that the item had been solved, but it prevents treating the old list alone
as proof of continuing openness.

The present search did not establish a previous publication of the
specific exclusion-closure argument, a payment of this particular reward,
or a current renewal of the reward. Establishing priority and applying
the 2007 monetary promise require review by the problem's authors.
The mathematical theorem here and those administrative questions should
be evaluated separately.

\medskip
\textbf{Preparation and reproducibility.}
This note was prepared with GPT-6 Astra assistance for Grisha Pochuev.
Independent reasoning passes checked the finite solver and the adaptive
sample argument. The supplied Python code requires only the standard
library. Run \texttt{python cotrain\_disjunctions.py --demo} for the example
and \texttt{python verify\_solver.py} for the implementation checks.
The source \LaTeX{} and exact verification results accompany the PDF.

\begin{thebibliography}{9}\small
\bibitem{bb07}
A. Blum and M.-F. Balcan.
\emph{Open Problems in Efficient Semi-Supervised PAC Learning}.
COLT 2007, pp.\ 622--624, especially Sections 3 and 3.2.
\url{https://home.ttic.edu/~avrim/Papers/openproblems.pdf}

\bibitem{bb10}
M.-F. Balcan and A. Blum.
\emph{A Discriminative Model for Semi-Supervised Learning}.
Journal of the ACM 57(3), 2010. Author version (2009).
\url{https://home.ttic.edu/~avrim/Papers/labeled_unlabeled-2009.pdf}

\bibitem{bbel13}
M.-F. Balcan, C. Berlind, S. Ehrlich, and Y. Liang.
\emph{Efficient Semi-supervised and Active Learning of Disjunctions}.
ICML 2013, PMLR 28(1), pp.\ 633--641.
\url{https://proceedings.mlr.press/v28/balcan13.html}

\bibitem{dmrt24}
I. Diakonikolas, M. Ma, L. Ren, and C. Tzamos.
\emph{Fast Co-Training under Weak Dependence via Stream-Based Active Learning}.
ICML 2024, PMLR 235, pp.\ 10841--10864.
\url{https://proceedings.mlr.press/v235/diakonikolas24b.html}
\end{thebibliography}
\end{document}

